\documentclass[11pt,a4paper]{article}
% Shared preamble for RuPhO English problem statements (match T1 2026 style).
\usepackage[margin=1in]{geometry}
\usepackage{amsmath,amssymb}
\usepackage{graphicx}
\usepackage{float}
\usepackage{enumitem}
\usepackage{microtype}
\usepackage{caption}
\usepackage{hyperref}

\captionsetup{font=small,labelfont=bf,skip=6pt}

\setlist[itemize]{leftmargin=1.2cm}
\setlist[enumerate]{leftmargin=1.2cm}

\newcommand{\problembox}[3]{%
  \par\addvspace{0.42em}%
  \noindent
  \setlength{\fboxsep}{7.5pt}%
  \setlength{\fboxrule}{0.95pt}%
  \fbox{%
    \begin{minipage}{\dimexpr\linewidth-2\fboxsep-2\fboxrule\relax}
      \textbf{#1}\hfill\textbf{#3}\par\smallskip
      #2
    \end{minipage}%
  }%
  \par\addvspace{0.32em}%
}

\title{\textbf{Josephson effect}}
\author{\normalsize X25 (2025) — English translation}
\date{}
\begin{document}
\maketitle
A superconductor is a substance through which current can flow at zero voltage without loss of energy. It turns out that if two superconductors are separated by a thin layer of non-superconducting material (for example, a normal metal with finite resistance or an insulator), then current can still flow through such a layer without resistance, provided that it does not exceed some critical value. Such a system, consisting of two superconductors separated by a non-superconducting layer, is called a Josephson junction, and the effect of current flowing through such a contact without loss of energy is called the Josephson effect. This effect was theoretically predicted by Josephson in 1962 and was soon discovered experimentally. The Josephson effect has several important practical applications, notably a magnetic field sensor and a voltage standard. The state of the superconductor is given by the complex number $\Delta = |\Delta|e^{i\theta}$; the square of the modulus of this number specifies the concentration of electrons in the superconducting state. The magnitude of the superconducting current through the contact is determined by the phase difference $\varphi=\theta_2-\theta_1$ between these two superconductors. In the simplest case (considered in this problem), the dependence has the form
$$
I_s = I_c \sin\varphi.
$$
Here $I_c$ is the critical current, depending on the material and dimensions of the contact.
\begin{figure}[H]
    \centering
    \includegraphics[width=0.72\textwidth]{T3_EN_assets/fig1.jpg}
    \caption{Fig. 1.}
\end{figure}
If a current greater than $I_c$ is passed through the contact, a voltage $V$ will appear across it, and the total current will be equal to the sum of the normal current $I_n=V/R$ and the superconducting current. In this case, the phase difference will change over time; its change is described by the equation
\[
2eV=\hbar\dot{\varphi}.
\]
Here $e$ is the electron charge and $\hbar$ is the reduced Planck constant. It is convenient to introduce the combination $\Phi_0=\pi\hbar/e$ (with the dimension of magnetic flux), so that
\[
V=\frac{\Phi_0}{2\pi}\,\dot{\varphi}.
\]
In addition, charges may accumulate at the edges of the contact; this can be taken into account by adding a capacitor in parallel with the contact. Thus, in our model the equivalent circuit of the contact consists of an ideal Josephson element (carrying current without resistance), a resistor $R$, and a capacitor $C$ connected in parallel, with the voltage across them given by the formula above.
\begin{figure}[H]
    \centering
    \includegraphics[width=0.72\textwidth]{T3_EN_assets/fig2.jpg}
    \caption{Fig. 2.}
\end{figure}
\section*{Part A. Contact without capacitance (5.3 points)}
In this part we will consider the behavior of a Josephson contact without taking into account capacitance, then the capacitor in the equivalent circuit can be replaced by an open circuit.
\problembox{A1}{Let at the initial moment of time $t = 0$ the current through the contact be equal to zero. At this moment, voltage $V$ was applied to the contact.  Determine the current versus time $I(t)$. Express your answer in terms of $\Phi_0$, $V$, $I_c$, $R$.}{0.50}
\problembox{A2}{Let a direct current $I > I_c$ pass through the contact.  Write down the differential equation for $\varphi(t)$.}{0.50}
\problembox{A3}{Differentiate the equation with respect to time. Eliminate the phases from the equation and obtain an expression for $\dot{V}(t)^2$ in terms of $V$, $I$ and system parameters.}{0.70}
\problembox{A4}{In the previous equation, go to the variable $u = 1/V$, express $\dot{u}^2$ in terms of $u$, $I$ and the system parameters.}{0.30}
\problembox{A5}{You should get an equation similar to the law of conservation of energy for oscillations of a harmonic oscillator. Determine the frequency of oscillations $\omega$, the value $u_0$ about which the oscillations occur, and the amplitude of oscillations $A$.}{0.80}
\problembox{A6}{Write down the dependence $V(t)$. Consider that at $t = 0$ the value of $V(t)$ has a minimum.}{0.20}
\problembox{A7}{Construct high-quality graphs of the dependence $V(t)$ for the cases $I = 1.1I_c$ and $I= 3 I_c$, indicate the coordinates of the characteristic points.}{0.60}
\problembox{A8}{Find the average (over time) voltage across the contact. Express your answer in terms of $R$, $I$, $I_c$.}{0.80}
Let no current initially flow through the contact. Let's connect a source to it and slowly increase the value of the current flowing through the contact from zero to a certain value $I< I_c$. The phase change during such a process will also be slow, so the normal component of the current through the contact can be neglected. This means there will be no energy losses either, and all the work of the source will go to increase the energy of the Josephson junction.
\problembox{A9}{Determine the dependence of the energy $E(\varphi)$ of the Josephson junction as a function of the phase difference $\varphi$ on it. The answer may also include $\Phi_0$ and $I_c$. Consider $E(0) = 0$.}{0.70}
\problembox{A10}{For a typical Josephson junction, the value of the critical current is $I_c = 0.5~mA$, the resistance is $R = 0.7 ~\Omega$. Fundamental constants $e = 1.602 \times 10^{-19}~C$, $\hbar = 1.055 \times 10^{-34}~J \cdot s$. Find the numerical value of the energy of the Josephson junction $E(\pi)$ and the frequency of voltage oscillations at constant current through the junction $I = 2 I_c$.}{0.20}
\section*{Part B. Contact with container (1.5 points)}
In this part we will consider the behavior of the contact taking into account its capacitance.
\problembox{B1}{Let a direct current $I$ flow into the contact (see the figure in the introduction to the problem). Write down a differential equation describing the change in phase difference $\varphi$. The answer may also include $\Phi_0$, $R$, $C$, $I_c$.}{0.30}
\problembox{B2}{Let the external current be zero and the phase difference $\varphi \ll 1$. In this case, $\varphi$ (and therefore other quantities characterizing the contact) will perform damped oscillations. Determine the vibration frequency $\omega_p$.}{0.60}
\problembox{B3}{Determine the quality factor of the oscillatory system from the previous paragraph $Q$. At what values of capacitance in the system are oscillations possible rather than damped motion?}{0.60}
\section*{Part C. Superconducting quantum interferometer (1.2 points)}
Consider a pair of Josephson contacts arranged as in the circuit in the figure. These contacts are formed by two superconductors separated by two layers of insulators so that a ring is formed. In the center of the ring there is a coil that creates a magnetic flux $\Phi$. Due to the magnetic field, the phase inside the conductor begins to depend on the coordinates, so the phase differences at the two contacts $a$ and $b$ are different from each other. If $\varphi_a$ and $\varphi_b$ are the phase differences at these contacts, they are related to the magnetic flux $\Phi$ by
\[
\varphi_a-\varphi_b = 2\pi\frac{\Phi}{\Phi_0}.
\]
Such a device is called a superconducting quantum interferometer (SQUID, superconducting quantum interference device).
\begin{figure}[H]
    \centering
    \includegraphics[width=0.72\textwidth]{T3_EN_assets/fig3.jpg}
    \caption{Fig. 3.}
\end{figure}
\problembox{C1}{Let the critical currents through the contacts be $I_{ac}$ and $I_{bc}$ respectively. Then the maximum superconducting current $I$ that can flow through such a contact will depend on the magnetic flux through the interferometer. Determine this maximum value $I_0$, as well as the cosines of the phase differences $\varphi_a$, $\varphi_b$ at which this maximum is achieved. Express your answer in terms of $I_{ac}$, $I_{bc}$, $\Phi$, $\Phi_0$. Your answers should not contain inverse trigonometric functions.}{1.20}
If the current through the interferometer exceeds the maximum value found in the previous paragraph, a voltage will arise across the interferometer, similar to the case of a single contact. By studying the behavior of this voltage, it is possible to measure magnetic fields with extremely high accuracy.
\vspace{1em}
\noindent\textit{Source:} \url{https://pho.rs/p/3448}
\end{document}
